\chapter{Interface utilisateur (React)}
\label{chap:frontend}

\section{Principes d'interface}
Le frontend, développé en \textbf{React} avec le bundler \textbf{Vite}, vise un usage par
des ingénieurs de maintenance, non par des spécialistes de la donnée. Il privilégie la
lisibilité immédiate (codes couleur sémantiques vert/jaune/orange/rouge), l'information en
temps réel (WebSocket) et des graphiques interactifs (\textbf{Recharts}). L'accès est
protégé par une page de connexion ; un intercepteur ajoute automatiquement le jeton JWT à
chaque requête et gère son expiration. Un \emph{error boundary} global garantit qu'une
erreur de rendu d'une page n'interrompt jamais toute l'application.

\section{Les pages fonctionnelles}
\subsection{Connexion}
Authentification par e-mail et mot de passe, avec gestion des rôles. La session est
conservée localement et expire automatiquement.

\subsection{Vue Globale}
Tableau de bord de la flotte : chaque machine est présentée par une carte affichant son
indice de santé (jauge), sa RUL estimée, sa tendance et son statut. Les KPIs agrégés
(machines saines, en alerte, critiques, indice moyen) sont en tête. Un clic sur une machine
déclenche sa surveillance temps réel et ouvre un panneau détaillé (indice de santé, RUL,
score d'anomalie, défaut détecté avec niveau de confiance).

\subsection{Monitoring}
Surveillance vibratoire en direct : forme d'onde, spectre et indicateurs (RMS, kurtosis,
crête) défilent en temps réel via le WebSocket, avec contrôles de lecture (play/pause,
vitesse) et sélection de la machine.

\subsection{IA Lab}
Centre d'analyse des modèles, organisé en onglets : \emph{Classification} et
\emph{Régression RUL} (graphiques et tableaux du benchmark), \emph{Fiabilité \& Drift}
(score de dérive par dataset, carte d'ensemble d'anomalie), \emph{Explicabilité SHAP}
(importance des indicateurs), \emph{Calibration} (diagramme de fiabilité, ECE) et
\emph{Versions Modèle} (historique des versions et \textbf{retour arrière} en un clic).
L'utilisateur peut y \textbf{basculer le modèle actif} d'un dataset à chaud.

\subsection{Pronostic \& RUL}
Visualisation de la trajectoire de dégradation (courbe d'indice de santé avec seuils
70/40/20\,\%) et compte à rebours de la durée de vie restante, avec comparaison RUL prédite
vs réelle pour les turbomoteurs CMAPSS.

\subsection{Maintenance}
Page opérationnelle centrale, enrichie de plusieurs panneaux : le \emph{journal d'alertes},
le \emph{copilot} conversationnel, le \emph{diagnostic par spectre externe} (import CSV/JSON),
la \emph{saisie d'événements de maintenance}, le panneau \emph{KPIs réels \& efficacité
prédictive} (MTBF, MTTR, disponibilité, ROI, précision/rappel/lead-time), le panneau
\emph{Ordres de travail} (boucle GMAO) et le \emph{réentraînement} piloté.

\subsection{Audit Trail}
Journal de traçabilité : chaque décision de l'IA et chaque alerte y est horodatée
(prédiction, machine, confiance), avec compteurs et filtres. Cette page matérialise
l'exigence \og rien n'est une boîte noire \fg.

\subsection{Configuration}
Réglage des seuils (anomalie, sévérité), des paramètres de roulement et des préférences de
la plateforme.

\section{Communication temps réel}
Le hook \texttt{useSimulationWS} encapsule la connexion WebSocket : ouverture, envoi des
commandes (play, pause, set\_machine), réception et bufferisation de l'historique, et
surtout \textbf{reconnexion automatique à back-off exponentiel} en cas de coupure. Les pages
consomment ce hook pour afficher des données vivantes sans recharger.

\section{Composants réutilisables}
Des composants d'interface (\texttt{HealthGauge} pour la jauge de santé, \texttt{MetricCard}
pour les indicateurs) assurent la cohérence visuelle. La charte graphique (fond sombre,
accents cyan/vert/ambre) a inspiré les couleurs de ce rapport.
